%% abstract.tex: abstract (EN)
%% -------------------------------------------------------------

\vspace{1cm}

\chapter*{Abstract}

The increasing deployment of autonomous robot teams in complex missions, from
search and rescue operations to remote industrial inspections, fundamentally
depends on these systems' ability to communicate reliably in environments where
network infrastructure is unavailable. In such situations, robots must
autonomously organize themselves into mobile ad-hoc networks, where each element
functions simultaneously as information source, recipient, and data relay,
enabling distant robots to communicate through multiple intermediate hops.

The main challenge in these networks is their intrinsically dynamic nature: as
robots move, communication links continuously establish and break, requiring the
routing system to constantly adapt routes to maintain active communication flows.
Previous work established solid foundations in this area, demonstrating functional
multi-hop routing on actual mobile robots. However, approaches at the data link
layer present limitations in terms of debugging complexity and integration with
standard networking tools.

This work presents an alternative architecture based on network layer routing,
integrated with the Linux kernel routing subsystem. The approach offers greater
robustness through proven kernel code, better observability through standard
tools, and explicit failure detection with bounded recovery time. The system uses
the connectivity information already distributed by the RA-TDMAs+ protocol to
compute routes automatically, recomputing them whenever it detects a change in the
network topology, and separates a UDP control plane for the topology beacons from
a TCP data plane for the application traffic.

The system was implemented and validated on a physical three-node testbed, in
which a mobile robot streams live video to a base station and an intermediate node
relays the traffic when the direct link fails. The evaluation measured routing
convergence after a link failure, round-trip latency, video throughput, and the
regularity of the stream, and confirmed that the nodes keep to their TDMA slots,
for both the periodic beacons and the continuous video, with the relay leaving
latency, throughput and timing essentially unchanged. The results show that
network layer routing is a robust and fully functional alternative for real
robotic applications, combining reliability, observability, and ease of
integration with existing systems.

\textbf{Keywords:} mobile ad-hoc networks, autonomous robotics, multi-hop
communication, dynamic routing, distributed systems
